\newpage
\sscp{\tsc{Sangir}}{C-09-01}
The following languages are all Austronesian,
belonging to the \tsc{Sangir} subgroup,
classified by \citet{Blust2019} as a member of \tsc{Philippine}.
\\

\noindent{\wg\st{3pt}\begin{tabularx}{\linewidth}
{lllll>{\raggedright\arraybackslash}X}
\lsptoprule
%Language	&	ISO	&	Term	&	Gloss	&	Etymon	&	Source	\\	\midrule
%\endfirsthead
Language	&	ISO	&	Term	&	Gloss	&	Etymon	&	Source	\\	\midrule
%\endhead
Bantik	&	bnq	&	kusey	&	cuscus	&	**kusay	&	\cite[89]{Sneddon1984}	\\
Ratahan	&	rth	&	kusey	&	cuscus	&	**kusay	&	\cite[89]{Sneddon1984}	\\
Sangir	&	sxn	&	kuse/	&	cuscus	&	**kusay	&	\cite[89]{Sneddon1984}	\\	\hhline{~}
				&			&	kusai	&		&		&	\hp{\,}	\\
Talaud	&	tld	&	usseʔa\footnotemark	&	cuscus	&	(**kusay)	&	\cite{BalaiBahasaSulawesiUtara2018}	\\
Talaud,	&	tld	&	uhheʔa	&	cuscus	&	(**kusay)	&	\hp{\,}	\\	\hhline{~}
Tirrawatta 	&		&		&		&		&	\multirow{-2}{=}{\raggedright \cite{BalaiBahasaSulawesiUtara2018}}	\\
\lspbottomrule
\end{tabularx}}
\footnotetext{ %The \it{uhheʔa} form is found in the Tirrawatta dialect of Talaud.
		A perusal of \citet{Sneddon1986} indicates that consonant gemination
		is found in Talaud,
		but the addition of root-final glottal stop appears to be irregular,
		more associated with other Sangir languages.
		Paragogic vowels are common in the region.
		Altogether it appears the Talaud forms are irregular from **kusay.}

%\noindent{\wg\st{2pt}\begin{xltabular}{\linewidth}
%{lllll>{\raggedright\arraybackslash}X}
%\lsptoprule
%%Language	&	ISO	&	Term	&	Gloss	&	Etymon	&	Source	\\	\midrule
%%\endfirsthead
%Language	&	ISO	&	Term	&	Gloss	&	Etymon	&	Source	\\	\midrule
%\endhead
%
%\lspbottomrule
%\end{xltabular}}